
Large language models have demonstrated proficiency in code generation, yet effective orchestration of multiple feedback sources (static analysis and execution testing) remains underexplored. This work investigates whether sequential single-source feedback routing provides computational efficiency advantages over simultaneous multi-source aggregation in LLM-based code generation. We implement cascade routing that applies static analysis (mypy --strict) before execution testing (pytest), conditionally gating execution when static errors exist. Using CodeLlama-7B on 35 dual-sensitive programming tasks from HumanEval, we measure static analysis error detection rates and token efficiency. Results show that mypy detects errors in 99.6\% of generation attempts (697/700 samples), substantially exceeding the predicted 30-40\% range. Mock simulation suggests token efficiency ratio of 0.733 (cascade vs. aggregation), though this awaits full experimental validation. One mechanism hypothesis (attention economy via sequential presentation) could not be tested due to implementation errors. Findings are scoped to dual-sensitive tasks with CodeLlama-7B base model on statically-typed Python code.
